%%% vim: set tw=72 %%%
%%% File:      introCTX_env.tex
%%% Author:    Joaquín Ataz-López et al
%%% Created:   March 2020
%%% File edited with: Emacs + AUCTeX - vim
%
% This file contains the preamble of the introduction to ConTeXt.  In
% the beginning I copied it from the source of the file about.pdf,
% included in the documentation on ConTeXt Standalone, and from there
% some things were changed in accordance with need.

\startenvironment introCTX_env.tex

    \mainlanguage[en] % Main language = English

    \setupbodyfont[modern] % The document's main (body) font

    % Page layout

    \setuplayout
        [
            width=middle,
            height=middle,
            topspace=2.5cm,
            backspace=3cm,
            header=1cm,
            headerdistance=0.5cm,
            footer=0.5cm
        ]

    % Use Roman numerals for the frontmatter:
    \defineconversionset[frontpart:pagenumber]
                        []
                        [romannumerals]

    % And when we start the bodymatter, reset page number to 1:
    \startsectionblockenvironment[bodypart]
        \setupuserpagenumber[number=1]
    \stopsectionblockenvironment

    \setupexternalfigures[directory=img, location=global] % Image directory

    \setupwhitespace[big] % Larger vertical space

    \setupitemize[indenting=none] % First line of paragraphs not indented

    \setupcaptions[headstyle=\bfx, style=\tfx] % Style of table and image titles

    \definecolor[maincolor] 
                [.6(orange)] % Title colour

    \definecolor[purple] 
                [.5(darkblue)] % Cover colour

    % Writing commands and code snippets
    \setuptype[color=darkmagenta]
    \setuptyping[color=darkmagenta, style={\switchtobodyfont[small, tt]}]

    % Labels
    \setuplabeltext
        [en]
        [
            chapter=Chapter~,
            appendix=Appendix~,
            part=Part~
        ]

    % Footnotes
    \setupnotation
        [footnote]
        [way=bypage, numberconversion=n, after={\blank[small]}]

    % Page number location on non-"/title/chapter/section/" pages:
    \setuppagenumbering[location={header, right}]
    % The cover page does not get a number (apparently \startXYZmakeup pages
    % don't get page numbers), but conceptually it is page 'i'.
    % Set the page number to 2 to take this unnumbered page into account.
    % (TODO: This is wrong when the cover page has not been included; fix!)
    \setupuserpagenumber[number=2]

    % Empty initial footers
    \setupfootertexts[]

    % CONFIGURATION OF PARTS, CHAPTERS, SECTIONS, ETC.

    % This snippet defines a footer with page numbers. A name is assigned
    % to it, and then, when section commands are configured, this name is
    % assigned to the first page of part and chapter.
    \definetext[ChapterFirstPage]
               [footer]
               [pagenumber]

    \setuphead
        [part]
        [
            placehead=yes,
            alternative=middle,
            conversion=I,
            style={\switchtobodyfont[40pt]\bf},
            color=purple,
            before={\ \godown[4cm]},
            header=high,
            footer=ChapterFirstPage,
        ]

    \setuphead
        [title, chapter]
        [
            style={\switchtobodyfont[30pt]\bf},
            align=flushright,
            before={\blank[5*big, force]},
            after={\blank[2*big]},
            color=maincolor,
            sectionsegments=2:2,
            commandbefore={\blank[small]},
            header=high,
            footer=ChapterFirstPage,
        ]

    % JD note: I hate seeing a section or subsection start right at the
    % bottom of a page, maybe leaving no lines of text after the title.
    % So add `\need' to the following setups.  
    % Why not \testpage[...]?  Good question.  See definition of \need below.
    % (And yes, the definition is plain TeX-y, not ConTeXt-y.  Mea culpa.)
    %
    % The \need 10\baselineskip below says "don't start a new section
    % within 10 lines of the bottom of the page".
    % Ditto (but smaller amounts) for subsection and subsubsection.
    \setuphead
        [section, subject]
        [
            style={\bfc},
            before={\blank[2*big]\need 10\baselineskip\relax },
            after=\blank,
            color=maincolor,
            sectionsegments=2:3,
        ]

    \setuphead
        [subsection, subsubject]
        [
            style={\bfb},
            before=\blank\need 8\baselineskip\relax,
            after=\blank,
            color=maincolor,
            sectionsegments=2:4,
        ]

    \define[1]\PlacePoint{#1.}

    \setuphead
        [subsubsection, subsubsubject]
        [
            sectionsegments=5:5,
            conversion=A,
            style=\bfa,
            before=\need 5\baselineskip\relax,
            color=black,
            numbercommand=,
        ]

    \setuphead
        [subsubsection]
        [numbercommand=\PlacePoint]


    % Structures defined by me

    \definedescription
        [description]
        [
            alternative=hanging,
            width=fit,
            distance=0.3cm,
            headstyle=\bf
        ]

    \definedescription
        [semitable]
        [
            alternative=left,
            width=fit,
            distance=1cm,
            headcolor=darkmagenta,
            headstyle=\tt,
        ]


    % ENVIRONMENTS

    % DoubleExample: For two column examples. It gave some errors when
    % defining it with \definestartstop (due to my lack of knowledge) and
    % in the end what I did was to define a TeX style command for opening
    % it, and another command for closing it.
    \def\startDoubleExample
    {
        \startframedtext
            [frame=off]
            \switchtobodyfont[small]
            \setuptyping[style=tt]
            \startcolumns
                [
                    n=2,
                    tolerance=verytolerant,
                    distance=0.5cm,
                    separator=rule,
                    rulecolor=darkmagenta
                ]
    }
    \def\stopDoubleExample{\stopcolumns\stopframedtext}

    % Small print
    \definestartstop[SmallPrint]
                    [
                        before={\startnarrower\switchtobodyfont[small,ss]},
                        after={\stopnarrower},
                    ]


    % Commands and abbreviations

    % An alternative "non-environment" for showing side-by-side examples of
    % code and results.  The code must be stored in a buffer.
    % The output is wrapped in a frame so that it stays together on a page.
    % Ironically, this uses some "beyond introductory"-level stuff.
    % Usage: \ShowAndExecute          --- \typebuffer, \getbuffer
    % Usage: \ShowAndExecute[Buf]     --- \typebuffer[Buf], \getbuffer[Buf]
    % Usage: \ShowAndExecute[Buf1][Buf2] --- \typebuffer[Buf1], \getbuffer[Buf2]
    %        (third case used when it you need to "cheat" for some good reason).

    \defineparagraphs[ShowAndExecute]
                     [n=2, rule=on, rulecolor=darkmagenta, 
                      distance=0.75cm, tolerance=verytolerant]
    \unprotect % Allow _ and @ in macro names
    \tolerant\def\ShowAndExecute [#1]#*[#2]
    {
        \startnarrow[middle=15pt]
            % Use a frame here to keep things together on one page.
            \startframedtext[offset=none,width=local,frame=off]
                \startparagraphs[ShowAndExecute]
                    \ifparameter#1\or
                        \typebuffer[#1]
                        \nextShowAndExecute
                        \switchtobodyfont[small]
                        % By default \\ means "start a new paragraph" in
                        % the paragraphs environment, but that is not
                        % desired in this case.  Deal with it.
                        \let\\=\crlf
                        \ifparameter#2\or
                            \getbuffer[#2]
                        \else
                            \getbuffer[#1]
                        \fi
                    \else
                        \typebuffer
                        \nextShowAndExecute
                        \let\\=\crlf			% See above
                        \switchtobodyfont[small]
                        \getbuffer
                    \fi
                \stopparagraphs
        \stopframedtext
            \stopnarrow
     }
     \protect


    % A (non-environment) command to show an *externally* typeset example
    % next to the corresponding code.
    % A frame is drawn around the typeset code to help visualize the page.
    % This was written to demonstrate footnotes; in MkIV footnote text
    % did not "migrate" out of a columns environment, but in LMTX it
    % does, so there had to be some updating of the examples in that
    % chapter.
    %
    % The size of the body font was chosen to make the text in the
    % typeset code more or less match that in ordinary side-by-side
    % examples.  This size (26pt at time of writing) is a function of the
    % font used by \ShowAndExecute and DoubleExample (small at time of
    % writing) as well as the width argument of \typesetbuffer (namely,
    % 0.45tw).
    %
    % This command requires 0, 1 or 2 buffers inside []:
    % 0: use the unnamed buffer for both sides
    % 1: use the argument buffer for both sides
    % 2: use #1 for the left (typed) side, #2 for the right (executed) side.
    %    (This is in case the typeset version needs a bit of "cheating"
    %    to make the output look good.)
    %
    % The optional third argument, which must be in {}, is used to
    % over-ride [...=...] arguments of \definepapersize.
    % See the macro itself for the default width and height.
    % Caution: typeset page is (currently) hard-coded to occupy 0.45tw,
    %          increasing the width of the typeset page using this parameter
    %          doesn't make the right side bigger, but it will make the
    %          text look smaller.
    %
    % TODO: this produces an overfull \hbox (8.51683pt).
    %       Where is that heppening, and how can it be fixed.
    %       The right edge is a bit closer (~=8pt) to the right margin
    %       than the left edge is to the left margin, but otherwise
    %       there is no apparent problem with the overfullness.    

    \defineparagraphs[ShowAndTypeset]
                     [n=2, rule=on, rulecolor=darkmagenta,
                      distance=0.75cm, tolerance=verytolerant]
    \unprotect % Allow _ and @ in macro names
    \tolerant\def\ShowAndTypeset[#1]#*[#2]#:#*#={%
        \ifparameter#1\or\def\_SAT_B{#1}\else\def\_SAT_B{{}}\fi
        \ifparameter#2\or\def\_SAT_BB{#2}\else\edef\_SAT_BB{\_SAT_B}\fi

        \startluacode
	  local preamble = [[
            \definepapersize[ShowAndTypeset]
                            [width=200mm, height=100mm, #3]
            \setuppapersize[ShowAndTypeset]
            % Margins: roughy 1" margins scaled down to 200/textwidth.
            % Top and bottom: leave just a bit of white space.
            \setuplayout[backspace=12mm, rightmargindistance=12mm,
                         edgedistance=0mm, margin=0mm, edge=0mm,
                         width=fit,
                         topspace=2mm, footer=2mm, footerdistance=2mm, 
                         height=fit]
            % This makes the text roughly the same as "small":
            \switchtobodyfont[26pt]
            % No need for a page number in these examples.
            \setuppagenumbering[location=]
            % Without this, footnotes are teeny-tiny
            \starttext
          ]]
          context.tobuffer("TypesetEnvironment", preamble)
        \stopluacode

        \startnarrow[middle=15pt]
            % Use a frame here to keep things together on one page.
            \startframedtext[offset=none,width=local,frame=off]
                \startparagraphs[ShowAndTypeset]
                    \ifparameter#1\or
                        \typebuffer[#1]
                        \nextShowAndTypeset
                        \switchtobodyfont[small]
                        \ifparameter#2\or
                            \typesetbuffer[TypesetEnvironment,#2]
                            		  [frame=on,width=0.45tw]
                        \else
                            \typesetbuffer[TypesetEnvironment,#1]
                              		  [frame=on,width=0.45tw]
                        \fi
                    \else
                        \typebuffer
                        \nextShowAndTypeset
                        \switchtobodyfont[small]
                        \typesetbuffer[TypesetEnvironment,{}]
                              	      [frame=on,width=0.45tw]
                    \fi
                \stopparagraphs
            \stopframedtext
        \stopnarrow
    }
    \protect


    % Assumption: Print the assumption icon in the margin
    \define\Conjecture
        {\inmargin[line=1]{\externalfigure[miniconjecture.jpg][width=1cm]}}

    % Doubt: Print the doubt icon in the margin
    \define\Doubt
        {\inmargin[line=1]{\externalfigure[minidoubt.jpg][width=1cm]}}

    % Example: Show in red and in small print the result of compiling an
    % earlier code snippet.
    \define[1]\example{\color[red]{\tfx#1}}

    % ConTeXt Standalone
    \def\suite-{\quotation{\ConTeXt\ Standalone}}

    % cmd: to be used in place of \tex when, in the source file, there is a
    % line break in an argument.
    % ALSO: it isn't verbatim, in that it evaluates the arguments.
    % JAL had this, but it uses the wrong font for \:
    %    {{\en\tt\color[darkmagenta]{\backslash #1}}}
    \define[1]\cmd
        {{\en\tt\color[darkmagenta]{\tex{}#1}}}

    % Since JAL uses \backslash a lot where he wants {\tt \}, define a
    % command to produce the correct character for (pseudo-)verbatim text
    % ("text back slash")
    % A horrible hack to get a \.  What is The True ConTeXt Way of doing this?
    \define\ttbs{\char92}
    \define\textbs{\char92}

    % PalClave:  tt and guillemets (angle quotes)
    % \define[1]\PalClave{«{\tt #1}»}

    % MyKey: tt and double inverted commas
    \define[1]\MyKey{\quotation{\tt #1}}

    % Partial chapter table of contents:
    \def\Separate#1{#1;}
    \def\TocChap
    {
        \startnarrower\switchtobodyfont[small, sl]
            {\bf Table of Contents:}
            \setuplist
                [section, subsection, subsubsection]
                [pagenumber=no, textcommand=\Separate]
            \setuplist[section][style=bold]
            \placecontent[alternative=d]
        \stopnarrower
    }

    % A command to forcibly eject this page if there is not the given
    % amount of space left on the current page.
    % In principle, in LMTX we should be able to \testpage[...], but
    % when I tried that in b03_Commands.tex context errored out.
    % In the fullness of time this should be investigated, but for
    % expediency here is a less sophisticated test that doesn't error out.
    % Usage:  \need <dimem>    E.g., \need 1cm
    \def\need{\afterassignment\zneed\dimen0=}
    \def\zneed{\relax
        \ifnum\pagetotal>\pagegoal  % Page may be full, but not yet gone
            % See if we can easily (0.3\pageshrink?) shrink the page to fit...
            % If so, bite the bullet & eject the page
            \dimen0=\pagetotal
            \advance\dimen0 by -0.3\pageshrink
            \ifnum\dimen0<\pagegoal
                \page % was \eject in plain TeX
            \else
                % V 1.23 change, see comments above
                \penalty-500
            \fi
        \else
            \vskip 0pt \par
            \ifnum\pagetotal>\pagegoal
                % At great risk, assume a \page is safe here, since all
                % we've added is glue to get past the end of the page
                \page % was \eject in plain TeX
            \fi
            \advance\dimen0 by \pagetotal
            \ifdim\dimen0>\pagegoal
                \page % was \vfill\eject in plain TeX
            \fi
        \fi
     }    

    % Index
    \defineregister[macro]
    \define[1]\PlaceMacro{\macro[#1]{\backslash #1}}

    % INTERACTIVITY

    \setupinteraction
      [
        state=start,
        color=darkblue,
        contrastcolor=darkblue,
        style=,
        title={A Not So Short Introduction to ConTeXt LMTX (Community Edition)},
        author={Joaquín Ataz-López et al},
      ]

    \setupurl [color=blue, style=\tt]

    \useurl[pragma]
           [http://www.pragma-ade.com/]

    \useurl[wiki]
           [https://www.contextgarden.net/]

    % Formerly [https://wiki.contextgarden.net/Symbols/utf8], but then 
    % Garulfoo changed things.
    \useurl[wikisymbols]
           [https://wiki.contextgarden.net/Characters_words_and_fonts/Emoji_and_symbols/utf8]

    % This must come after \setupinteraction.  Ask me not why.
    \placebookmarks[chapter, section, subsection, subsubsection]

    % TWO INSTRUCTIONS ADDED BY J. TO ELIMINATE MOST HYPHENATION AND GET
    % BETTER INTERWORD SPACING AS WELL

    % JD note: I have seen cases where using 'hz' causes anomalous line
    % breaking when paragraphs are being formed.  Specifically, increasing one
    % inter-word space from "normal size" to "sentence ending size" (which
    % causes the line to get slightly longer), made the algorithm add yet
    % another word to that line.  Which is at least counterintuitive.
    % TODO: I wonder whether we want it or whether we are better off without it.
    \setupalign[hz, lesshyphenated, verytolerant, stretch]

    % JD: I disagree with this one... I believe the advice "don't use two
    % spaces after a period" is coming from the word processor world, which
    % doesn't have the same smarts as TeX about having *slightly* larger
    % spaces following full stops; TeX doesn't use "two spaces" anyway.
    % I personally find slightly larger spaces makes reading easier.
    % So until such time as I am shouted down...
    % \setuplanguage[en][spacing=packed]% This avoids an extra space
    %% after full stops (periods) which is no longer 'de rigeur' in printed
    %% material in English.

    \useURL[SourceCode]
           [https://codeberg.org/anss-ce/ANSS-CE][][]

\stopenvironment

%%% Local Variables:
%%% mode: ConTeXt
%%% eval: (auto-fill-mode 1)
%%% TeX-master: "introCTX_eng.tex"
%%% coding: utf-8-unix
%%% End:
